% =====================================================================================================
% Chương 4, PHẦN B (mục 4.6–4.11) — bản viết lại (2026-10-12).
% Nối ngay sau 08-chuong4-phanA.tex trong cùng \input{noi-dung/08-chuong4}; không có \chapter mới.
% Ghi chú cho người ghép: 09-chuong4-phanB-NOTES.md (nhãn, nguồn bảng/hình, bằng chứng thiếu, chỗ lệch outline).
% Mỗi dòng "% EV:" trỏ tới EVIDENCE-REGISTER.md (A01…K52), EVIDENCE-REGISTER-ADDENDUM-1.md (N-01…N-30,
% L1a-01…09, mục 0.x), EVIDENCE-REGISTER-ADDENDUM-2-CH3.md (X3-xx) hoặc EVIDENCE-REGISTER-ADDENDUM-3-CH4B.md
% (N4B-01…N4B-14, tạo cùng phần này).
% Không dùng gói mới: hình vẽ bằng TikZ thường; danh sách dùng enumitem như chương 3.
% =====================================================================================================

% -----------------------------------------------------------------------------------------------------
\section{Vai trò của trần và thứ tự ưu tiên chọn phương án}
\label{muc:c4-vaitro}

Điều phối không trần và luật C trần 30 giây dùng cùng một tệp cấu hình và cùng bộ chọn ưu tiên ổn định, chỉ
khác nhau ở trần giờ trả. Chênh lệch giữa hai cấu hình ở các mục~\ref{muc:c4-ketqua}--\ref{muc:c4-congdon} vì
vậy là chênh lệch đi kèm việc thêm trần khi giữ nguyên bộ chọn, đo trên hai tập khách được phục vụ khác nhau.
Bộ chọn này vốn xếp tiêu chí ổn định trước chi phí kế hoạch (tiêu chí 3 và 4 ở Bảng~\ref{bang:rangbuoc}), nên
còn một câu hỏi mở: đánh đổi của luật C có còn khi bộ chọn ưu tiên chi phí hay không.
% EV: A08; H01; NEW:12 (tệp U và C-30 chỉ khác /policies/0/limits/1/hardLimit); C04; A14; X3-07; H02

Bằng chứng hiện có là một phép sàng lọc trên tập tắc đường cố định. Ở mỗi lượt quyết định đã ghi, bộ chọn được
chạy lại trên chính các phương án đã qua phép kiểm của cấu hình đó, sau khi bỏ tiêu chí 3; các tiêu chí khác giữ
nguyên. Lựa chọn mới chỉ được so tại trạng thái mà lần chạy đã đi qua; mô phỏng không chạy tiếp từ lựa chọn đó.
Với thứ tự tiêu chí gốc, phép chạy lại cho đúng lựa chọn đã ghi ở cả 62.837 lượt quyết định của năm cấu hình.
Bảng~\ref{bang:c4-sangloc} so hai cấu hình; mỗi cột tính trên các trạng thái của chính cấu hình đó.
% EV: N4B-01 (r17_wp14r_check.py:1-19, 32; WP17:3-22); H03 (PV:2-6; WP17:3,6,10,15,20); H10

\begin{table}[htbp]
\centering
\caption{Phép sàng lọc bỏ tiêu chí ổn định của bộ chọn, trên tập tắc đường cố định (16 kịch bản mỗi cấu hình,
bộ chạy cũ). Mỗi cột dùng các lượt quyết định đã ghi của chính cấu hình đó, nên hai cột là hai tập trạng thái
khác nhau. Chi phí kế hoạch là thời lượng còn lại của các tuyến dự kiến, không phải chi phí vận hành.}
\label{bang:c4-sangloc}
\small
\begin{tabular}{@{}p{10.6cm}rr@{}}
\toprule
\textbf{Chỉ số} & \textbf{U} & \textbf{C-30}\\
\midrule
Lượt quyết định có yêu cầu mới & 1.714 & 1.715\\
Lượt đổi lựa chọn khi bỏ tiêu chí ổn định & 445 (25,96\%) & 104 (6,06\%)\\
Trong đó, lượt mà lựa chọn mới dời thêm ít nhất một khách đang có & 426 & 81\\
% EV: H04 (445/1.714; 104/1.715; số yêu cầu được nhận không đổi); H09 (426/445; 81/104; DER17:15,18)
\midrule
\multicolumn{3}{@{}l}{\textit{Chi phí kế hoạch giảm (\%), ba cách tính}}\\
Trung bình trên các lượt đổi lựa chọn & 7,79 & 8,00\\
Trung bình trên mọi lượt có yêu cầu mới & 2,02 & 0,49\\
Tổng phần giảm chia tổng chi phí kế hoạch đã chọn, trên mọi lượt quyết định & 0,3069 & 0,0794\\
% EV: H05 (NEW:6; PV:2,5); K15; A16; I09
\midrule
\multicolumn{3}{@{}l}{\textit{Tổng dời giờ trả do điều phối tăng thêm trên đội xe, ở các lượt đổi lựa chọn (s)}}\\
Trung vị & 90,8 & 12,5\\
Lớn nhất & 390,8 & 35,5\\
% EV: H06 (WP17:8,22; PV:2,5); K17; CASE:1-18
\bottomrule
\end{tabular}
\end{table}

Cách đọc bảng: với điều phối không trần, trong 1.714 lượt có yêu cầu mới, bỏ tiêu chí ổn định làm lựa chọn đổi
ở 445 lượt; ở các lượt này chi phí kế hoạch giảm trung bình 7,79\%, còn tổng dời giờ trả do điều phối trên đội
xe tăng thêm, trung vị 90,8 giây.
% EV: H04; H05; H06

Số yêu cầu được nhận không đổi ở lượt nào: bộ chọn ưu tiên chi phí chỉ chọn cách chở khác cho cùng các yêu
cầu. Ba cách tính chi phí trả lời ba câu hỏi khác nhau nên không thay được cho nhau; chi phí kế hoạch lại được
tính lặp qua các lượt, nên không cộng thành mức tiết kiệm vận hành. Tỉ lệ lượt đổi lựa chọn là 25,96\% khi không
trần, 18,95\% (325 trong 1.715) với trần 120 giây, 10,08\% (173 trong 1.716) với trần 60 giây và 6,06\% với trần
30 giây. Mỗi tỉ lệ đo trên trạng thái của chính cấu hình, nên đây không phải tác động của trần trên cùng một
trạng thái; kết quả chỉ gợi ý rằng dưới phép kiểm của trần 30 giây, bộ chọn có ít cơ hội đổi hơn trên các trạng
thái đã kiểm. Lượt có tổng dời tăng thêm lớn nhất dưới trần 30 giây, 35,5 giây, chia cho hai khách, mỗi khách
17,77 giây.
% EV: H04; H05; K15; A16; I09; K16; H06; K17 (CASE:1-18: hai khách, mỗi khách 17.770 ms)

Phép sàng lọc không cho biết số lượt phục vụ, trễ hay đi chung trên quỹ đạo, và không chuyển sang tập kiểm tra
chính, vì tập tắc đường cố định dùng thời gian đi không đổi và bộ chạy cũ không có luật phục hồi. Phép thử đổi thứ
tự ưu tiên của bộ chọn trên 34 kịch bản chính chưa được chạy; bằng chứng hiện có là phép sàng lọc trên trạng thái
cố định ở trên, vốn không cho biết kết quả phục vụ và đúng giờ trên quỹ đạo. Vì vậy, khóa luận chưa tách được phần
của trần và phần của bộ chọn trong độ tin cậy đã đo; chương kết luận đặt phép thử này lên đầu các hướng phát triển.
% EV: H10; NEW:13 (khác thiết lập bộ chạy); B17; I12; H12 (chưa chạy); REG mục L (L02–L05 chưa có kết quả); H02; K18

% -----------------------------------------------------------------------------------------------------
\section{Lời hứa theo dự báo}
\label{muc:c4-dubao}

Trên 21 kịch bản ngày thường có dự báo, số khách của luật C trần 30 giây tới muộn hơn lời hứa đầu quá một phút
giảm từ 68 xuống 2, trong 1.471 khách được phục vụ (gộp theo khách). Tính trung bình theo kịch bản, tỉ lệ này
giảm từ 6,01\% xuống 0,17\%.
% EV: F04 (6,01% → 0,17% trung bình theo ô; gộp 2/1.471); F17 (2 so với 68); RP1:39-40; K24

Lời hứa theo dự báo ở đây là cách tính của mục~\ref{muc:dubao}, phát lại trên các lần chạy đã ghi: xe nào chở ai
giữ nguyên, chỉ lời hứa đầu đổi, và trễ được so với lời hứa đầu của chính cách hứa đó. Đối chứng thứ hai là
\emph{đệm cùng giá}: mỗi lời hứa đầu được kéo dài thêm một tỉ lệ cố định của thời gian còn lại tới giờ trả, tỉ lệ
chọn riêng cho từng khung giờ sao cho tổng phần kéo dài trên các khách của luật C bằng tổng phần kéo dài của dự
báo. Bảng~\ref{bang:c4-dubao} so ba cách hứa này theo khung giờ.
% EV: F03 (chỉ vế phát lại); F01; N4B-03 (forecast_family_a.py:15-17: lời hứa đệm = P0 + k_w(P0 − t0); FA:2)

\begin{table}[htbp]
\centering
\caption{Khách của luật C trần 30 giây tới muộn hơn lời hứa đầu quá 60 giây, theo ba cách hứa, trên 21 kịch bản
ngày thường có dự báo (1.471 khách được phục vụ). Mọi số là trung bình theo kịch bản, tính bằng \%; dòng cuối là
tỉ lệ khách tới sớm hơn lời hứa quá 60 giây.}
\label{bang:c4-dubao}
\small
\begin{tabular}{@{}lrrr@{}}
\toprule
\textbf{Nhóm kịch bản} & \textbf{Không dự báo} & \textbf{Đệm cùng giá} & \textbf{Dự báo}\\
\midrule
Cả 21 kịch bản & 6,01 & 1,94 & 0,17\\
Khung 7 giờ (8 kịch bản) & 14,44 & 4,24 & 0,27\\
Khung 8 giờ (8 kịch bản) & 1,35 & 0,86 & 0,18\\
Khung 17 giờ (5 kịch bản) & 0 & 0 & 0\\
% EV: F04; N4B-03 (FR:40-45; FA:8-11: cột đệm 1,94 / 4,24 / 0,86 / 0)
\midrule
Tới sớm quá 60 s, cả 21 kịch bản & 0 & chưa đo & 0,08\\
% EV: N4B-03 (FR:46; FA:8: E naive 0, disp 0,0008; đệm không tính E)
\bottomrule
\end{tabular}
\end{table}

Cách đọc bảng: ở khung 7 giờ, trung bình mỗi kịch bản có 14,44\% khách tới muộn hơn lời hứa đầu quá một phút
khi không dự báo, 4,24\% khi đệm cùng giá và 0,27\% khi có dự báo.
% EV: F04; N4B-03

Chênh lệch dự báo trừ không dự báo trên 21 kịch bản là $-5{,}84$ điểm phần trăm. Khoảng tin cậy 95\% tính bằng
cách lấy mẫu lặp (bootstrap) theo 8 cụm ngày và khung giờ là từ $-8{,}78$ đến $-3{,}23$; một cách tính thận trọng
hơn theo phân phối $t$ trên chính 8 cụm này cho khoảng từ $-10{,}54$ đến $-1{,}15$. Mức giảm có ở cả ba buổi sáng
và vẫn còn khi bỏ riêng từng ngày. So với đệm cùng giá, dự báo thấp hơn 1,77 điểm phần trăm, nhưng khoảng theo
phân phối $t$ chạy từ $-3{,}62$ đến $+0{,}07$, tức chạm 0: với ít cụm như vậy, lợi thế so với cách đệm này còn
mong manh. Ngày 12/11 là một ngày nguồn của mọi dự báo; bỏ ngày này thì chỉ còn hai buổi sáng có dự báo và khoảng
theo phân phối $t$ rộng ra từ $-18{,}9$ đến $+6{,}4$.
% EV: F05 (FA:32-42; FR:51-60); N4B-04 (FA:36,38-39: đệm [−2,83; −0,68], từng buổi sáng, bỏ từng ngày); F11 (FR:94-99); I05

Lợi ích này không riêng của luật C. Trên cùng 21 kịch bản, luật hạn chót và luật trần độ dời khách thấy, đều ở
mức 30 giây, cũng xuống 0,17\%; với điều phối không trần, tỉ lệ giảm từ 9,72\% xuống 3,90\% (trung bình theo kịch
bản). Riêng ở khung 17 giờ, tỉ lệ 6,76\% của điều phối không trần không đổi khi có dự báo, trong khi luật C ở khung
này đã là 0\% từ trước. Dự báo và trần nhắm vào hai phần khác nhau của độ lệch giờ trả.
% EV: F08 (FR:65,79-86; FA:3,6,13,23); K37; FA:11 (C-30 khung 17 giờ 0)

Giá của dự báo là lời hứa dài hơn. Trên 21 kịch bản, lời hứa trả dài thêm trung bình 13,7 giây, ở khung 7 giờ là
27,8 giây, và lời hứa đón dài thêm 3,2 giây (trung bình theo kịch bản). Trên cả 34 kịch bản
(Bảng~\ref{bang:c4-tran120}), số lượt trễ quá một phút của luật C giảm từ 76 xuống 10 với lời hứa dài thêm trung
bình 7,37 giây; số này nhỏ hơn vì nó chia cho cả 2.338 khách được phục vụ, kể cả khách của 13 kịch bản cuối tuần,
nơi không có dự báo và phần dài thêm bằng 0. Dự báo cũng có chỗ chưa khớp: phép kiểm khả thi vẫn dùng giờ không dự
báo, nên 34 trong 1.471 khách có giờ đón hiển thị ở lời hứa đầu muộn hơn giờ đón muộn nhất (mục~\ref{muc:dubao}).
% EV: F06 (FR:67-69; FA:8-9); C09; N-16; N4B-10 (trung bình gộp gồm khách cuối tuần, phần dài thêm 0); F02; N-28; F07

Kết quả này phải đọc cùng một giới hạn của mô phỏng. Tắc đường là một hệ số chung cho cả thành phố, nên khi kéo
dãn lời hứa theo hệ số thật thay vì hệ số dự báo, giờ trả khớp trong vòng 2 giây với mọi khách có phần dời do điều
phối bằng 0: 1.452 trên 1.452 khách ở 21 kịch bản và 2.305 trên 2.305 khách ở 34 kịch bản. Thế giới mô phỏng vì vậy
gần như tất định, và 0,17\% là mức mà một dự báo theo giờ trong ngày đạt được trong thế giới ấy, không phải hiệu
quả ngoài đời.
% EV: F10 (FR:77; RX:8; FA:28); I02; B06

Sổ theo từng khách (công thức~\eqref{eq:so}) cho biết dự báo giải thích phần nào của độ trễ. Với 68 khách của luật
C tới muộn hơn lời hứa đầu không dự báo quá một phút ở 21 kịch bản ngày thường, 6.217 giây trễ gồm 82 giây tính cho
điều phối, 5.353 giây được hồ sơ lịch sử giải thích và 783 giây còn lại (Hình~\ref{hinh:c4-so}). Hình đặt cạnh đó
chiều tuyết của mục~\ref{muc:c4-tuyet}, nơi tỉ lệ gần như đảo ngược. Đây là phép tính theo mô hình: phần thứ hai phụ
thuộc dự báo đang dùng, và phần còn lại không phải phần ``không thể biết trước''.
% EV: F09 (FR:72-75; FA:30; RX:3); N4B-05 (chiều tuyết); K04; K19

\begin{figure}[htbp]
\centering
\begin{tikzpicture}[x=1cm,y=1cm,font=\footnotesize,
  nhan/.style={font=\scriptsize,inner sep=1.5pt}]
% Thanh trên: ngày thường. Độ dài 11 cm = 6.217,4 s; ranh giới 0,1445 và 9,6145 cm (FA:30).
\draw[fill=black!75] (0,1.6) rectangle (0.1445,2.2);
\draw[fill=black!20] (0.1445,1.6) rectangle (9.6145,2.2);
\draw[fill=white] (9.6145,1.6) rectangle (11,2.2);
\node at (4.88,1.9) {5.353 s (86,1\%)};
\node[align=center,nhan] at (10.307,1.9) {783 s\\(12,6\%)};
\draw (0.072,2.2)--(0.072,2.45);
\node[nhan,anchor=south west] at (0.0,2.45) {82 s (1,3\%)};
\node[anchor=east,align=right] at (-0.15,1.9) {Ngày thường, 21 kịch bản\\68 khách, 6.217 s};
% Thanh dưới: chiều tuyết. Độ dài 11 cm = 3.445,9 s; ranh giới 0,0273 và 0,8302 cm (NEW:5).
\draw[fill=black!75] (0,0) rectangle (0.0273,0.6);
\draw[fill=black!20] (0.0273,0) rectangle (0.8302,0.6);
\draw[fill=white] (0.8302,0) rectangle (11,0.6);
\node at (5.915,0.3) {3.185,9 s (92,5\%)};
\draw (0.014,0)--(0.014,-0.2);
\node[nhan,anchor=north west] at (-0.02,-0.2) {8,6 s};
\draw (0.43,0)--(0.43,-0.62);
\node[nhan,anchor=north west] at (0.39,-0.62) {251,5 s (7,3\%)};
\node[anchor=east,align=right] at (-0.15,0.3) {Chiều tuyết, 8 kịch bản\\41 khách, 3.445,9 s};
% Chú giải
\draw[fill=black!75] (0,-1.55) rectangle (0.3,-1.3);
\node[anchor=west] at (0.35,-1.425) {phần tính cho điều phối};
\draw[fill=black!20] (3.7,-1.55) rectangle (4.0,-1.3);
\node[anchor=west] at (4.05,-1.425) {phần được hồ sơ lịch sử giải thích};
\draw[fill=white] (8.6,-1.55) rectangle (8.9,-1.3);
\node[anchor=west] at (8.95,-1.425) {phần còn lại};
\end{tikzpicture}
\caption{Ba phần của số giây trễ theo sổ theo từng khách (công thức~\eqref{eq:so}), với luật C trần 30 giây và các
khách tới muộn hơn lời hứa đầu không dự báo quá 60 giây: 68 khách của 21 kịch bản ngày thường và 41 khách của 8
kịch bản chiều tuyết 15/11/2018. Mỗi thanh dài bằng 100\% số giây trễ của nhóm đó. Phần được hồ sơ lịch sử giải
thích phụ thuộc dự báo đang dùng; phần còn lại là phần dự báo này chưa giải thích được, không phải phần không thể
biết trước.}
\label{hinh:c4-so}
\end{figure}
% EV (ngày thường): F09 (FA:30: 81,7 + 5.352,6 + 783,0 = 6.217,4 s; RX:3: 1,3 / 86,1 / 12,6%)
% EV (chiều tuyết): F13 (8,6 trên 3.445,9 s); N4B-05 (NEW:5: 8,566 / 251,496 / 3.445,933 s; RP1:24-26: 7,3%; SNOW:10: 92,5%; FR:104,186)
% Tọa độ: x = giây / tổng × 11 cm (tính tay; xem NOTES mục 4).

\textbf{Đưa dự báo vào bộ điều phối.} Một cách khác là cho chính bộ điều phối dùng thời gian đi theo dự báo
(mục~\ref{muc:dubao}). Cách này được thử với luật C và luật hạn chót, đều ở mức 30 giây, trên 16 kịch bản khung 7
và 8 giờ. Trong 32 lần chạy, 3 lần dừng giữa chừng vì FleetPy, chạy theo hệ số đo, tính xe tới một điểm đón muộn hơn
giờ đón muộn nhất 1,67 hoặc 2,69 giây. Luật C không phải phục hồi lần nào ở 14 lần chạy hoàn tất. Luật hạn chót vẫn
phải phục hồi ở 47,6\% lượt quyết định, so với 62,8\% khi không đưa dự báo vào (trung bình trên 15 kịch bản chạy
xong); tỉ lệ khách bị chính nó phá lời hứa ở khung 7 giờ giảm từ 39,0\% xuống 11,0\%. Con số phá lời hứa ở
Bảng~\ref{bang:c4-nghiavu} vì vậy phụ thuộc cách hứa. Do có lần chạy dừng, cách triển khai được chọn là tính dự báo
ở lớp lời hứa như các kết quả trên; kết luận này chỉ áp cho cách đưa dự báo vào đã thử.
% EV: F12 (FR:117-146; RX:9-15); X3-29; N4B-07 (FR:29,34,120-121; RX:9,11,13-15); C14 (lưu ý: phụ thuộc cách hứa)

\textbf{Học máy.} Khóa luận cũng thử thay phần kéo dài chung bằng một mô hình học máy ước lượng độ trễ của từng
khách. Ba mô hình được cố định, không dò tham số: cây tăng cường độ dốc (gradient boosting) ước lượng phân vị 90 của
độ trễ từ các đặc trưng có lúc hứa, như giờ trong ngày, thời gian còn lại tới giờ trả và số khách khác trên xe; cùng
mô hình đó nhưng bỏ số khách trên xe; và hồi quy phân vị tuyến tính. Các mô hình học trên tập buổi sáng 12--13/11 và
được kiểm trên ngày 15 và 16/11; trên tập tắc đường cố định, học trên ngày 12/11 và kiểm trên ngày 13/11. Mỗi mô hình
được so với một cách đệm đều có cùng độ dài thêm trung bình, mức chính là 30 giây. Một mô hình chỉ được coi là tốt
hơn khi thỏa đồng thời bốn điều kiện: tỉ lệ trễ quá một phút thấp hơn; độ dài thêm trung bình không quá 1 giây so
với đối chứng; phân vị 90 của độ dài thêm không quá 15 giây so với đối chứng; tỉ lệ tới sớm quá một phút tăng không
quá 1 điểm phần trăm.
% EV: G01; N4B-08 (MLS:26-35, 41-61, 63-80, 91-95); G02; G03

Kết quả không cho cải thiện nào thỏa đủ bốn điều kiện. Với luật C có dự báo, mô hình không giảm thêm tỉ lệ trễ ở cả
hai ngày kiểm. Với điều phối không trần, mức giảm là 0,69 điểm phần trăm ở ngày 15/11 (145 khách) và 0 ở ngày
16/11, nhưng lời hứa dài thêm quá mức cho phép. Với ghép tự do trên tập tắc đường cố định, tỉ lệ trễ giảm từ
12,89\% xuống 11,65\% trên 807 khách của ngày 13/11, nhưng 7,1\% khách tới sớm quá một phút, so với 0\% khi đệm
đều. Bỏ số khách trên xe không đổi kết quả nào, và không mô hình nào đạt cùng lúc tỉ lệ trễ và độ chính xác hai
phía của luật C có dự báo. Kết luận chỉ áp cho các mô hình cố định với đặc trưng đơn giản đã thử.
% EV: G05; G06; G07; G08; G09; G10 (câu được phép, HO:346-347); N4B-08 (ML:6,7,14,15,21: n và các điều kiện không thỏa); K40

% -----------------------------------------------------------------------------------------------------
\section{So với cách đệm lời hứa}
\label{muc:c4-dem}

Khi có dự báo, cộng thêm 30 giây vào mọi lời hứa của điều phối không trần cũng cho 2.328 lượt không trễ quá một
phút trên 3.672 yêu cầu, bằng luật C trần 30 giây, và cho 2.366 lượt không trễ quá hai phút, so với 2.338 của luật
C. Đổi lại, lời hứa dài thêm trung bình 37,2 giây thay vì 7,4 giây.
% EV: C10; N-16; N-18; K06

Ở mục này, \emph{đệm 30 giây} nghĩa là cộng đúng 30 giây vào lời hứa đầu theo dự báo của mọi khách. Phép đệm cũng
là phát lại trên các lần chạy đã ghi: quyết định điều phối, tập khách được phục vụ và giờ trả thực tế của điều phối
không trần giữ nguyên, chỉ lời hứa đổi. Vì vậy cách đệm không phải một đối chứng ghép cặp với luật C: hai bên phục
vụ hai tập khách khác nhau, chung 2.207 khách. Bảng~\ref{bang:c4-dem} xếp các cách hứa thành ba khối theo độ dài
thêm trung bình của lời hứa, để so các cách có lời hứa dài gần bằng nhau.
% EV: C10 (lưu ý: phát lại, không phải đối chứng ghép cặp); ADDENDUM-1 phần đầu (quy tắc đệm D = D0 + k); C04; R25 §6

\begin{table}[htbp]
\centering
\caption{Các cách hứa trên tập kiểm tra chính khi có dự báo (34 kịch bản, 3.672 yêu cầu mỗi cấu hình; dự báo chỉ có
ở 21 kịch bản ngày thường). Số đếm gộp; trễ tính trong số khách được phục vụ, so với lời hứa của chính cách hứa đó;
hai cột ``không trễ quá'' tính trên 3.672 yêu cầu; độ dài thêm là trung bình gộp theo khách được phục vụ.}
\label{bang:c4-dem}
\small
\begin{tabular}{@{}p{4.4cm}rrrrr@{}}
\toprule
\textbf{Cách hứa} & \textbf{Phục vụ} & \textbf{Trễ quá 60 s} & \textbf{Không trễ quá 60 s} & \textbf{Không trễ quá 120 s}
& \textbf{Dài thêm (s)}\\
\midrule
\multicolumn{6}{@{}l}{\textit{Không đệm (dài thêm khoảng 7 giây)}}\\
U, dự báo & 2.434 & 137 & 2.297 & 2.347 & 7,18\\
C-30, dự báo & 2.338 & 10 & 2.328 & 2.338 & 7,37\\
C-120, dự báo & 2.387 & 65 & 2.322 & 2.379 & 7,27\\
% EV: N-16; N-03; C09; N-06; L1a-03
\midrule
\multicolumn{6}{@{}l}{\textit{Đệm 30 giây (dài thêm khoảng 37 giây)}}\\
U, dự báo, đệm 30 s & 2.434 & 106 & 2.328 & 2.366 & 37,18\\
C-30, dự báo, đệm 30 s & 2.338 & 2 & 2.336 & 2.338 & 37,37\\
% EV: C10; N-02; N-16
\midrule
\multicolumn{6}{@{}l}{\textit{Đệm 60 giây (dài thêm khoảng 67 giây)}}\\
U, dự báo, đệm 60 s & 2.434 & 87 & 2.347 & 2.386 & 67,18\\
% EV: C11; N-16
\bottomrule
\end{tabular}
\end{table}

Cách đọc bảng: khi đệm 30 giây trên điều phối không trần, 2.434 yêu cầu được phục vụ và 106 lượt tới muộn hơn lời
hứa đã đệm quá 60 giây, nên 2.328 lượt không trễ quá 60 giây trên 3.672 yêu cầu.
% EV: C10; N-16

\textbf{Vai trò của từng thành phần.} Bốn điều sau tách phần mà dự báo, đệm và trần mỗi thứ mang lại.

Thứ nhất, so ở độ dài thêm gần bằng nhau, luật C hơn ở ngưỡng một phút và kém ở ngưỡng hai phút. Khi chưa đệm,
luật C có 2.328 so với 2.297 lượt không trễ quá một phút và 2.338 so với 2.347 lượt không trễ quá hai phút. Khi cả
hai cùng đệm 30 giây, các cặp số là 2.336 so với 2.328 và 2.338 so với 2.366. Cách đệm chỉ bằng luật C ở ngưỡng
một phút khi lời hứa của nó dài thêm khoảng năm lần (37,18 so với 7,37 giây).
% EV: N-17; N-18 (37,1833 / 7,3671 = 5,05); N-16

Thứ hai, lời hứa dài thì ít trễ nhưng nhiều khách tới sớm (Bảng~\ref{bang:c4-haiphia}). Với luật C có dự báo, 4
trong 2.338 khách tới sớm quá một phút; với cách đệm 30 giây là 44 trong 2.434 khách. Số lượt tới trong khoảng
$\pm$60 giây quanh lời hứa là 2.324 với luật C và 2.284 với cách đệm, trên 3.672 yêu cầu. Đệm 60 giây nâng số lượt
không trễ quá hai phút lên 2.386, nhưng 1.033 trong 2.434 khách tới sớm quá một phút.
% EV: C10; N-04; C11; N4B-09 (ADD1:14-26,84-89)

Thứ ba, phần đuôi khác nhau. Với cách đệm 30 giây, 68 khách vẫn trễ quá hai phút, 11 khách quá năm phút và 1 khách
quá mười phút; luật C có dự báo không có khách nào trễ quá hai phút. Lợi thế 28 lượt của cách đệm ở ngưỡng hai phút
đến cùng 96 lượt phục vụ nhiều hơn, không phải từ ít khách trễ hơn.
% EV: N-01; N-19

\begin{table}[htbp]
\centering
\caption{Hai phía của độ lệch so với lời hứa khi có dự báo, trên tập kiểm tra chính (34 kịch bản). Các cột trễ và
tới sớm đếm trong số khách được phục vụ (2.434 với U, 2.338 với C-30, 2.387 với C-120); cột ``trong $\pm$60 s''
đếm số lượt tới cách lời hứa không quá 60 giây, tính trên 3.672 yêu cầu.}
\label{bang:c4-haiphia}
\small
\begin{tabular}{@{}p{4.4cm}rrrrr@{}}
\toprule
\textbf{Cách hứa} & \textbf{Tới sớm quá 60 s} & \textbf{Trong $\pm$60 s} & \textbf{Trễ quá 120 s} &
\textbf{Trễ quá 300 s} & \textbf{Trễ quá 600 s}\\
\midrule
U, dự báo & 8 & 2.289 & 87 & 17 & 2\\
C-30, dự báo & 4 & 2.324 & 0 & 0 & 0\\
C-120, dự báo & 2 & 2.320 & 8 & 0 & 0\\
U, dự báo, đệm 30 s & 44 & 2.284 & 68 & 11 & 1\\
C-30, dự báo, đệm 30 s & 44 & 2.292 & 0 & 0 & 0\\
U, dự báo, đệm 60 s & 1.033 & 1.314 & 48 & 8 & 1\\
\bottomrule
\end{tabular}
\end{table}
% EV: N-01; N-02; N-03; N-04; N-06; C11; N4B-09 (ADD1:14-16,19-20,24,84-89; A:7,9: late600)

Thứ tư, cách đệm không giới hạn phần dời do điều phối của từng khách. Vì đệm không đổi quyết định điều phối, cách
đệm giữ nguyên 174 khách bị điều phối dời giờ trả, 12 khách bị dời từ hai lần và một khách bị dời tổng cộng 590,1
giây; dưới luật C trần 30 giây là 33 khách, không ai bị dời hai lần, nhiều nhất 29,2 giây
(Bảng~\ref{bang:c4-doidieuphoi}). Giới hạn này không cần dự báo đúng: nó vẫn giữ ở 13 kịch bản cuối tuần không có dự
báo và ở chiều tuyết (mục~\ref{muc:c4-dieuphoi}).
% EV: C20; N-23; C10 (lưu ý: đệm là phát lại trên quỹ đạo của U); N-28; N-26; ADDENDUM-1 mục 0.4

So với luật hạn chót, khác biệt nằm ở nghĩa vụ nào giữ được. Khi có dự báo, luật hạn chót 30 giây cũng phục vụ
2.338 lượt, có 2.328 lượt không trễ quá một phút và 2.338 lượt không trễ quá hai phút. Nhưng nó phải phục hồi ở
41,3\% lượt quyết định (trung bình theo kịch bản) và phá lời hứa của chính nó ở 28 trong 34 kịch bản
(Bảng~\ref{bang:c4-nghiavu}); luật C không phải phục hồi lần nào trên 34 kịch bản.
% EV: N-16 (lưu ý: M⁺-30 + dự báo 2.338 / 10 / 2.328 / 2.338, ADD1:29); C13; N-21; X3-16

Trần còn là một núm vặn có thể kết hợp với đệm. Trần 120 giây có dự báo cho 2.379 lượt không trễ quá hai phút,
nhiều hơn cách đệm 30 giây (2.366), với lời hứa dài thêm 7,27 giây, khoảng một phần năm của cách đệm, và phục vụ ít
hơn điều phối không trần 47 lượt. Đây là một điểm đo thêm trên đường đánh đổi, trên một tập khách khác, không phải
mức trần tốt nhất.
% EV: N-06; N-18 (0,1956; 47); N-10; K36

\textbf{Khi nào chọn trần, khi nào chọn đệm.} Nếu hãng xe chỉ tính số lượt không trễ quá hai phút trên mỗi yêu cầu,
và chấp nhận lời hứa dài hơn cùng nhiều khách tới sớm hơn, cách đệm trên điều phối không trần là một lựa chọn cạnh
tranh và phục vụ thêm 96 lượt. Nếu hãng xe cần một giới hạn cho từng khách về phần dời do ghép thêm, giữ được cả khi
dự báo sai hay không có, cùng với lời hứa ngắn, thì luật C cho điều đó; đổi lại là 96 lượt phục vụ và 242 khách đi
chung trên 3.672 yêu cầu. Hai cách cũng dùng chung được, như hàng luật C có dự báo và đệm 30 giây. Khóa luận không
có bằng chứng khách hay hãng xe coi trọng lời hứa ngắn hoặc giới hạn theo từng khách đến mức nào, nên chọn cách nào
là quyết định của hãng xe, không phải kết luận của thực nghiệm. Mọi so sánh ở mục này đều giữ bộ chọn ưu tiên ổn định
(mục~\ref{muc:c4-vaitro}).
% EV: N-19; C02; N-20; N-02; D20; I07; I10; K38; K51; H02

% -----------------------------------------------------------------------------------------------------
\section{Khi dự báo không đủ: chiều tuyết 15/11/2018}
\label{muc:c4-tuyet}

Trên 601 khách được phục vụ của chiều tuyết 15/11/2018, dự báo lịch sử chỉ giảm số khách tới muộn hơn lời hứa đầu
quá một phút từ 41 xuống 34 (gộp theo khách: từ 6,82\% xuống 5,66\%; trung bình theo kịch bản: từ 7,19\% xuống
5,94\%).
% EV: F13; K25; B20

Tập này gồm 8 kịch bản khung 14 và 15 giờ, chạy với luật C trần 30 giây và lời hứa không dự báo; bộ chạy của tập
chưa được kiểm lại (Bảng~\ref{bang:c4-taptu}). Theo sổ, trong 3.445,9 giây trễ của 41 khách, chỉ 8,6 giây tính cho
điều phối và 251,5 giây được hồ sơ lịch sử giải thích; 3.185,9 giây (92,5\%) là phần hồ sơ này chưa giải thích được
(Hình~\ref{hinh:c4-so}).
% EV: B20; F13; N4B-05; N4B-06 (rider_nowcast_replay.py:6: 8 lần chạy C-30 không dự báo); phần A Bảng tập thử (bộ chạy chưa kiểm lại)

Để xem thông tin cùng ngày giúp được bao nhiêu, khóa luận phát lại thêm vài cách hứa trên cùng các lần chạy
(Bảng~\ref{bang:c4-tuyet}). Các cách ``trộn'' lấy trung bình nhân giữa mức tăng theo hồ sơ lịch sử và xu hướng tăng
gần nhất của hệ số tắc đường trong ngày. Khi chỉ dùng hệ số đã có trong khung giờ, số khách trễ còn 21. Thêm hệ số
cùng ngày từ 6 giờ sáng thì còn 19, nhưng bộ điều phối không được giao phần dữ liệu buổi sáng này, và một quy tắc
chuyển sang cách hứa này khi tắc đường cao bất thường cho cùng con số với ngưỡng chọn sau khi xem dữ liệu. Khi xu
hướng chỉ dùng các khung đã qua thì còn 29, dù mức hiện tại vẫn lấy đúng từ bảng thời gian đi. Một cách đệm có cùng
độ dài thêm, được biết trước giá của ngày tuyết, cho 27.
% EV: F14; K20; K21; N4B-06 (factor_nowcast_screen.py:12-30; factor_nowcast_screen_v2.py:6-13; RP1:7-20; NEW:5)

\begin{table}[htbp]
\centering
\caption{Phát lại các cách hứa trên 8 kịch bản chiều tuyết 15/11/2018 (luật C trần 30 giây, 601 khách được phục
vụ). Cột ``khách trễ'' đếm gộp số khách tới muộn hơn lời hứa đầu quá 60 giây; cột cuối là độ dài thêm trung bình
theo kịch bản. Mọi cách hứa đều tính lại sau khi chạy, trên các ngày đã xem.}
\label{bang:c4-tuyet}
\small
\begin{tabular}{@{}p{3.6cm}p{7.0cm}rr@{}}
\toprule
\textbf{Cách hứa} & \textbf{Thông tin dùng lúc hứa} & \textbf{Khách trễ} & \textbf{Dài thêm (s)}\\
\midrule
Không dự báo & Thời gian đi hiện tại & 41 & 0\\
Dự báo lịch sử & Hồ sơ tắc đường của các ngày nguồn & 34 & 2,0\\
Trộn trong khung & Thêm hệ số cùng ngày đã có trong khung giờ & 21 & 8,5\\
Trộn có buổi sáng & Thêm hệ số cùng ngày từ 6 giờ, không được giao cho khung này & 19 & 9,5\\
Trộn trễ một khung & Như hàng trên, xu hướng chỉ dùng các khung đã qua & 29 & 9,0\\
Đệm cùng giá & Tỉ lệ đệm lấy theo hàng ``trộn trong khung'' trên chính ngày tuyết & 27 & 8,3\\
Hệ số thật & Hệ số sau lúc hứa, không có ở thời điểm hứa & 0 & 18,8\\
\bottomrule
\end{tabular}
\end{table}
% EV: F14 (RP1:7-20; NEW:5); N4B-06 (định nghĩa các cách hứa; cột dài thêm là trung bình theo kịch bản, RP1:6, rider_nowcast_replay.py:96-104)

Cách đọc bảng: khi trộn hồ sơ lịch sử với xu hướng trong khung giờ, 21 trong 601 khách tới muộn hơn lời hứa quá một
phút, với lời hứa dài thêm trung bình 8,5 giây mỗi kịch bản.
% EV: F14; N4B-06

Theo sổ, mô hình có dữ liệu buổi sáng giải thích thêm 886,9 giây (25,7\%) ngoài dự báo lịch sử, còn 2.299,0 giây
(66,7\%) vẫn chưa được giải thích bằng các cách đã thử. Đây là một sự kiện đã được xem trước, nên chỉ là một ví dụ.
Trong giai đoạn 9/2018--12/2019 chỉ có 5 đợt tuyết từ 25,4 mm trở lên, kể cả đợt này, nên một phép kiểm dự báo cho
ngày có sự kiện cần dữ liệu đúng thời điểm và nhiều sự kiện độc lập hơn.
% EV: F15; K19; I11; B26

% -----------------------------------------------------------------------------------------------------
\section{Kiểm tra tính hợp lệ của kết quả}
\label{muc:c4-hople}

Trên 34 kịch bản của tập kiểm tra chính, không khách nào vượt trần: tổng dời do điều phối lớn nhất của một khách là
29,2 giây với trần 30 giây và 118,9 giây với trần 120 giây, và không có lượt quyết định nào phải phục hồi (0 trong
20.101 và 0 trong 20.435 lượt).
% EV: C20; N-07; X3-16; N-21

Các giới hạn tính toán của bộ sinh không cắt kết quả trên tập này. Ở cả bốn cấu hình được kiểm (điều phối không
trần, luật C trần 30 và 120 giây, luật hạn chót 30 giây), giới hạn 100 phương án mỗi xe chưa lần nào phải bỏ phương
án hợp lệ, và giới hạn khối lượng khám phá chưa lần nào để sót nhánh chưa xét. Giới hạn 2 yêu cầu mới mỗi lượt chỉ
một lần phải để một yêu cầu đang chờ sang lượt sau, trong 34 lần chạy của mỗi cấu hình. Vì vậy chênh lệch giữa các
cấu hình phản ánh tác động của cả gói cam kết, không phải của giới hạn tính toán. Các bộ đếm này không biến bộ sinh
thành phép vét cạn: nó chỉ xét các cách chèn yêu cầu mới vào kế hoạch đang có, không đảo điểm cũ và không chuyển
khách sang xe khác (mục~\ref{muc:gioihan}).
% EV: N-25 (ADD1:119-122; A:7-9); N4B-11 (GEN:54-95, 298-345; CandidateModels.cs:100-118; R15-audit.py:76-79); K34; X3-03; X3-32; X3-39; A24

Phần đầu của tiêu chí 3 không phân biệt phương án trên tập này. Ở mọi lượt quyết định có ghi tiêu chí này (3.503
lượt khi không trần, 3.506 với trần 30 giây, 3.504 với trần 120 giây), mức dùng trần lớn nhất đều bằng 100\%, vì đổi
xe, đổi điểm đón và đổi điểm trả có trần 0. Tiêu chí 3 vì vậy chỉ tác dụng qua tổng phần dời theo từng chiều.
% EV: A15; N-25; A07; X3-07

Tệp cấu hình của điều phối không trần và của luật C trần 30 giây chỉ khác nhau ở một giá trị là trần giờ trả. Tệp
chính sách cam kết của điều phối không trần và của luật C ở mức 30, 60 và 120 giây giống nhau từng byte giữa tập kiểm
tra chính và tập tắc đường cố định; hai tập chỉ khác ở thiết lập bộ chạy, gồm luật phục hồi, cách ghi lần đón muộn
và phiên bản chính sách.
% EV: A08; NEW:11-13; N4B-12

Các số của tập kiểm tra chính được ít nhất hai bộ đọc viết độc lập tính lại từ bản ghi thô và khớp trong phạm vi đã
kiểm, gồm các trường của từng khách ở 7.159 khách được phục vụ của điều phối không trần và luật C trần 30, 120 giây;
trên tập này, mọi khách được nhận đều được chở tới nơi. Mục~\ref{muc:kiemtra} nêu thêm phép kiểm 382.006 bản ghi sửa
và 13 hạng mục được tính lại. Trên đúng mã nguồn đã dựng bộ chạy, khi chạy lại ngày 12/10/2026 trên một bản sao mới
với kiểu xuống dòng LF, 1.024 trong 1.025 kiểm thử tự động đạt; kiểm thử còn lại đòi cây làm việc có tệp chưa được
ghi nhận và đạt khi chạy riêng với một tệp như vậy.
% EV: N4B-14 (ADD1:66-78; ADDENDUM-1 phần đầu); ADD1:5-8 (khớp đã nhận = đã chở); K41; A26; C19; N4B-13 (TESTRUN-45f8026-2026-10-12.md; log dòng 43-105; log dirtycheck)

Các phép kiểm trên cho biết con số khớp bản ghi và cài đặt giữ đúng trần trên các lần chạy đã kiểm. Chúng không cho
biết mô phỏng giống thực tế đến đâu; mục~\ref{muc:c4-thaoluan} bàn giới hạn đó.
% EV: I01; K41

% -----------------------------------------------------------------------------------------------------
\section{Thảo luận và hạn chế}
\label{muc:c4-thaoluan}

Với một hãng xe, các kết quả của chương là một bảng các điểm đánh đổi, không phải một khuyến nghị. Trên tập kiểm
tra chính, so với điều phối không trần, luật C trần 30 giây đổi 96 lượt phục vụ và 242 khách đi chung lấy ba điều:
phần dời do điều phối của mỗi khách không quá 30 giây, ít lượt trễ hơn ở cả bốn ngưỡng đã đo, và không lượt quyết
định nào phải phục hồi. Trần 120 giây là một điểm ở giữa; cách đệm trên điều phối không trần là một lựa chọn khác,
phục vụ nhiều hơn với lời hứa dài hơn. Khóa luận không đo giá trị của các điều này với khách hay với hãng xe.
% EV: C02; N-20; C20; N-24; X3-16; N-06; C10; D20; I10; K38

Một phép tính giả định cho thấy cách đọc bảng này. Nếu mỗi lượt trễ quá một phút so với lời hứa đầu làm hãng xe tốn
một khoản $k$ và mỗi chuyến mang lại biên lợi nhuận $m$, thì với 96 lượt phục vụ ít hơn và 127 lượt trễ ít hơn, luật
C chỉ có lợi hơn điều phối không trần khi $k > 0{,}756\,m$ (tức $96/127$), nếu mọi yếu tố khác bằng 0. Hai nhóm
khách khác nhau nên đây không phải phép đổi một lượt lấy một lượt, và khóa luận không có giá trị thật nào của $k$ hay
$m$. Một giả thuyết chưa kiểm: ở mức tải cao hơn, mỗi lần ghép thêm chạm tới nhiều khách đang trên xe hơn, nên trần
có thể chặn nhiều phương án hơn và cái giá về lượt phục vụ có thể lớn hơn; thực nghiệm ở đây không cho biết chiều
hay độ lớn của thay đổi đó.
% EV: C29 (R15:231-235; 96/127 = 0,756); K35; C28; I10; câu cuối là giả thuyết, không có dòng sổ (RD LR-34); B12

Các hạn chế dưới đây chia theo bốn nhóm; mỗi ý nêu điều chưa làm và kết luận chịu ảnh hưởng.

\textbf{Mô hình mô phỏng.}
\begin{itemize}[leftmargin=1.2cm,itemsep=2pt,topsep=3pt]
\item Mọi kết quả đến từ một bộ mô phỏng, FleetPy, nối với RideBound. Vì vậy chúng không cho biết hiệu quả khi vận
hành thật.
\item Tắc đường là một hệ số chung cho cả thành phố theo khung 15 phút, với mức nền lấy lúc 00:00 cho mọi kịch bản.
Thế giới vì vậy gần như tất định, nên mức trễ còn lại sau dự báo là mức của thế giới này, không chuyển sang mạng có
tắc đường khác nhau trên từng đoạn đường.
\item Thời gian dừng đón, trả bằng 0 giây (kiểm trên 3 lần chạy mẫu), trong khi nghiên cứu điển hình của FleetPy dùng
30 giây mỗi khách~\cite{engelhardt2026}; đồ thị là bao đóng 96 nút, không phải mạng đường phố. Độ dài chuyến và độ
lệch giờ trả ở đây vì vậy chưa tính các yếu tố này.
\item Mô phỏng không có phản ứng của khách (không hủy chuyến, không từ chối lời mời) và không cấu hình nào cho chuyển
khách đã nhận sang xe khác. Vì vậy số lượt phục vụ không tính khách bỏ đi vì lời hứa, và chưa biết cái giá của trần
đổi thế nào khi được gán lại xe.
\end{itemize}
% EV: I01; B01; B06; B07; I02; F10; B05 (SERV:2-4; FLEET:423-424); J18; B03; K48; I07; X3-39

\textbf{Dữ liệu và thiết kế thực nghiệm.}
\begin{itemize}[leftmargin=1.2cm,itemsep=2pt,topsep=3pt]
\item Tập kiểm tra chính có 34 kịch bản thuộc 5 ngày, phụ thuộc nhau, mỗi kịch bản chạy một lần. Các phép đếm theo
kịch bản vì vậy là mô tả, không phải kiểm định, và không bảo đảm cho ngày khác.
\item Ngưỡng phân tích, cách chia tập và mức 120 giây được chọn sau khi xem dữ liệu phát triển; chiều tuyết là một
sự kiện đã xem. Mọi kết quả là thăm dò.
\item Tải thấp: trung bình 0,535 khách trên xe kể cả lúc xe rỗi với điều phối không trần, so với khoảng 1,33 trong
nghiên cứu điển hình của FleetPy với 180 xe và 21.400 yêu cầu; hai thiết lập khác nhau. Các chuyến được chọn cũng
ngắn hơn tập đủ điều kiện (trung vị đi thẳng 165 so với 196 giây, ở cả 49 kịch bản). Kết quả không suy ra cho tải
cao hay chuyến dài hơn.
\end{itemize}
% EV: B09; I05; B11; I06; B10 (lưu ý); I11; B12 (LOAD:2-3; FLEET:433-434,447-449); J18; B13 (SEL6:2-4)

\textbf{Tích hợp.}
\begin{itemize}[leftmargin=1.2cm,itemsep=2pt,topsep=3pt]
\item Lời hứa theo dự báo được tính sau khi chạy, trong khi phép kiểm khả thi vẫn dùng giờ không dự báo (34 trong
1.471 khách có giờ đón hiển thị muộn hơn giờ đón muộn nhất), và cách đưa dự báo vào bộ điều phối đã thử làm 3 trong
32 lần chạy dừng. Muốn triển khai phải sửa phép kiểm giờ đón trước.
\item Một bộ chạy dựng lại từ mã nguồn có thể khác bản lưu trữ; bộ chạy của chiều tuyết chưa được kiểm lại. Lần chạy
mới phải cho đúng kết quả cũ trên vài kịch bản trước khi đem so với các số ở đây.
\end{itemize}
% EV: F07; F12; I13; R18 B24 (qua REG A28); phần A Bảng tập thử (chiều tuyết: chưa kiểm lại)

\textbf{Phạm vi kết luận.}
\begin{itemize}[leftmargin=1.2cm,itemsep=2pt,topsep=3pt]
\item Bộ chọn ưu tiên chi phí chưa được chạy trên quỹ đạo (mục~\ref{muc:c4-vaitro}), và luật mỗi lần một trần chưa
được cài đặt (mục~\ref{muc:c4-congdon}). Vì vậy chưa tách được phần của trần với phần của bộ chọn, và chưa biết vế
cộng dồn có đổi kết quả hay không.
\item Luật hạn chót chỉ là phần hạn chót, ở 30--60 giây thay vì 900 giây của Schulz và Pfeiffer~\cite{schulz2026};
hạn chót nới kiểu Hörl và Zwick~\cite{horl2022} chưa được chạy. Các phép so vì vậy không nói về phương pháp của các
tác giả đó.
\item Mức 30 giây không phải ngưỡng cảm nhận của khách hay chuẩn ngành, và không có dữ liệu về mức khách sẵn lòng trả
cho lời hứa ổn định hơn. Vì vậy không mức trần nào được chọn là tốt nhất.
\end{itemize}
% EV: H02; H12; K18; E09; K08; J10; J11; K27; D20; K52; K36
